\documentclass[10.5pt,a4paper]{article}
\usepackage[T1]{fontenc}
\usepackage{lmodern}
\usepackage[margin=2cm]{geometry}
\usepackage{array,longtable,booktabs}
\usepackage{xcolor}
\usepackage{enumitem}
\usepackage{amsmath}
\usepackage[hidelinks]{hyperref}
\usepackage{microtype}

\definecolor{head}{HTML}{1F3A5F}
\definecolor{qcol}{HTML}{7A2E0E}
\newcommand{\var}[1]{\texttt{#1}}
\newcolumntype{L}[1]{>{\raggedright\arraybackslash}p{#1}}
\setlist{itemsep=2pt,topsep=3pt}
\setlength{\parindent}{0pt}
\setlength{\parskip}{4pt}
\renewcommand{\arraystretch}{1.25}

% two-column decision table: left = what I did, right = why and the trade-off
\newenvironment{decisions}{%
  \begin{longtable}{L{0.36\textwidth} L{0.58\textwidth}}
  \toprule
  \textcolor{head}{\textbf{What I did}} & \textcolor{head}{\textbf{Why, and what it costs (trade-off)}}\\
  \midrule\endhead}{%
  \end{longtable}}
\newcommand{\row}[2]{#1 & #2\\ \midrule}

% interview questions
\newenvironment{qa}{\par\noindent\textcolor{qcol}{\textbf{Likely interview questions}}\begin{description}[font=\normalfont\itshape,leftmargin=0pt,style=nextline]}{\end{description}}
\newcommand{\q}[2]{\item[``#1'']#2}

\title{\textbf{AI Adoption and the Performance of Indian Banks}\\[3pt]
\large Research decisions, trade-offs and reasoning --- interview preparation}
\author{Research project under Prof.\ Paras Sachdeva, DMS, IIT Delhi}
\date{Status as of 26 November 2026 (results in \texttt{results/summary.md}, every estimate in \texttt{results/spec\_log.md})}

\begin{document}
\maketitle

\section*{How to use this document}
Each section has a side-by-side table: the decision on the left, the reasoning and its cost on the right, followed by
the questions an interviewer is most likely to ask, with short answers in the first person. Numbers are copied from the
logged results; units are stated as in the project (ROE/ROA in percentage points, pp; stock return in percent, \%).

\section{The question and the one-paragraph pitch}
\textbf{Research question.} Does AI adoption improve the financial performance and market value of Indian banks, and is
the effect weaker for public sector (PSU) banks?

\textbf{Pitch.} I built a three-level design: (1) sector-level descriptive trends, (2) a bank-by-year panel of 25
listed banks, FY2011--FY2026, estimated with two-way fixed effects and staggered difference-in-differences, and (3) an
event study of 30 usable AI announcements across 14 listed banks, NBFCs and fintechs. I measured AI adoption six ways
(annual-report text, verified first deployment, event counts, a fintech index, an AI-capability index, earnings-call
text) because no single public measure of bank AI use exists in India. The analysis plan was written before
estimation, every deviation was logged with a date and reason before re-running, and every estimate --- including the
nulls --- is in the specification log.

\textbf{Headline findings (stated carefully).}
\begin{itemize}
\item Annual-report AI intensity (\var{ai\_kw}) has no detectable average effect on ROE (+0.13 pp per SD, $p=0.86$) or
  stock returns ($-1.72\%$, $p=0.70$).
\item The AI~$\times$~PSU interaction on ROE is positive (+4.39 pp, $p=0.007$, wild-bootstrap $p=0.007$) --- the
  \emph{opposite} sign to the hypothesis --- but a post-hoc check with PSU-specific year effects cuts it to +1.66 pp
  ($p=0.34$), so it mostly reflects the PSU post-AQR recovery.
\item Staggered DiD gives large negative ATTs (ROE $-23.9$ pp; stock return $-37.2\%$), but with significant ROE
  pre-trends and a small, unbalanced never-treated group, so I do not read them as causal. Post-hoc checks that remove
  the PSU-dominated control group shrink the ROE effect to $-3.3$ pp ($p=0.57$, private banks only) and $-2.5$ pp
  (Sun--Abraham with PSU-specific year effects); the negative stock-return effect remains ($-36.3\%$ and $-15.8\%$).
\item AI announcements produce no average abnormal return; PSU announcements show $-2.70\%$ CAR[$-5,+5$] against the
  Nifty 500 (Kolari--Pynn\"onen $p=0.035$), not robust to the Nifty Bank benchmark.
\end{itemize}

\section{Overall design}
\begin{decisions}
\row{Three phases: sector level (descriptive), firm panel (main), event study (market reaction).}
{Each level answers a different question: long-run trend (sector), within-bank performance effects (panel), and how
investors price the news on the day (event study). Triangulation matters because each level is weak alone. Cost: more
work and more places for inconsistency, so every phase uses the same fiscal-year convention, firm IDs and event file.}
\row{Phase 1 (sector) is \emph{descriptive only}.}
{There are only 14--15 group-year observations per level. Any OLS at that size is underpowered and cannot separate AI
from the cleanup of bad loans (AQR), demonetisation or COVID. I say this upfront rather than over-claim.}
\row{Wrote a fixed analysis spec before estimation; every change logged in \var{spec\_changes.md} with date and reason
\emph{before} re-running; every estimate (incl.\ nulls) in \var{spec\_log.md}.}
{Protects against specification searching (``garden of forking paths''). With six proxies, four outcomes and several
interactions, something will be significant by chance; logging everything lets a reader count the tests. Cost: I
cannot quietly drop an awkward result --- which is the point.}
\row{Free, public data only (RBI DBIE, BSE/NSE filings, Yahoo prices, annual reports).}
{Fully reproducible and auditable; anyone can re-download. Cost: no CMIE Prowess, so some variables (market cap,
employee/tech data) had to be constructed, and some proxies (patents) turned out to have no usable data.}
\end{decisions}
\begin{qa}
\q{Why not just run one regression?}{Because the three questions are different. The panel tells me about accounting
and annual returns; the event study isolates the short-run price reaction to the news itself, where identification is
cleaner; the sector view shows the big picture but is not causal.}
\q{How do you know you didn't p-hack?}{The spec was fixed first, every change was logged before re-running, and the
spec log holds all 192 estimated rows, including the null ones. The post-hoc analyses (excluding Yes Bank, the
AI~$\times$~PSU check, the DiD parallel-trends diagnostics, the same-firm event-overlap rule) are labelled post-hoc
everywhere and the pre-specified estimates stay the headline.}
\end{qa}

\section{Samples and time frame}
\begin{decisions}
\row{Panel: 25 listed banks (12 PSU, 13 private), FY2011--FY2026; IDFC First only from FY2019 (post-merger).
Unbalanced panel, 392 bank-years.}
{Banks share one regulator (RBI), one accounting framework and one data source (RBI bank-wise tables), so ROE and
controls are comparable. NBFCs and fintechs report under different rules and are excluded from the panel. Cost: 25
clusters is few for inference (handled with the wild bootstrap).}
\row{Event study: 14 fixed firms (5 private banks, 4 PSU banks, 2 NBFCs, 3 fintechs).}
{The market reaction can be measured for any listed firm, so the event study can include the firms most active in AI
(Paytm, PB Fintech, Bajaj Finance). Razorpay and PhonePe are unlisted, so qualitative only.}
\row{Fiscal year = 1 April to 31 March (FY2017 = Apr 2016--Mar 2017).}
{Matches Indian bank reporting and RBI tables, so accounting and stock-return years line up.}
\row{Accounting outcomes end FY2025; stock returns run to FY2026.}
{RBI's FY2026 bank-wise tables were not yet published. I did not fill them from other sources, to keep one consistent
source. Logged as a spec change.}
\end{decisions}
\begin{qa}
\q{Why exclude NBFCs and fintechs from the panel?}{Different regulation, capital rules and accounting; their ROE is not
comparable with banks', and RBI's bank-wise tables do not cover them. They stay in the event study, where only
returns are needed.}
\q{Isn't 25 banks too small?}{It is most of the listed Indian banking sector by assets, but it is few clusters, so I
report wild cluster bootstrap $p$-values and treat borderline results cautiously.}
\end{qa}

\section{Outcomes and how they were built}
\begin{decisions}
\row{Primary outcomes: ROE = PAT / average equity $\times100$; annual stock total return (1 Apr--31 Mar) from
adjusted close. Secondary: ROA and Tobin's Q.}
{ROE is what bank management and analysts target; stock return captures what the market expects. Using average equity
avoids distortion from large capital infusions (common for PSU banks). Adjusted close already includes dividends and
splits, so no separate dividend add-on (avoids double counting).}
\row{Market cap built as (paid-up capital / face value) $\times$ FY-end price; face-value splits taken from split
records; cross-checked against Yahoo's current market cap (must be within 5\%).}
{No free historical market-cap series exists. Cost: bank-years where RBI's paid-up capital includes preference capital
(government-held, for some PSU banks and Kotak) would overstate shares, so those 39 bank-years are set to NA rather than
guessed.}
\row{Tobin's Q = (market cap + total liabilities) / total assets.}
{Standard approximation for banks (book liabilities $\approx$ market value). Cost: noisy; reported as secondary.}
\row{Suspicious data set to NA, never patched (e.g.\ KVB FY2012 return, unreliable Yahoo data before May 2011).}
{Rule: never invent data; every missing value is logged. A wrong number is worse than a missing one.}
\end{decisions}
\begin{qa}
\q{Why is ROE primary and not ROA?}{ROE is the return to shareholders and the bank's headline target; ROA is reported
as secondary. Results are reported for all four outcomes so the choice does not hide anything.}
\q{How did you get market capitalisation without Prowess?}{Shares outstanding from RBI paid-up capital divided by face
value, adjusted for splits, times the FY-end price, validated against Yahoo within 5\% for 5 banks. Where preference
capital contaminated the share count, I set the value to missing.}
\end{qa}

\section{Measuring AI adoption: six proxies}
\subsection*{Why six}
There is no administrative data on bank AI use in India. Each proxy captures a different facet (talk, deployment,
announcements, digital infrastructure, organisational capability, management attention) and each has a different
bias. If results agree across proxies, that is evidence; if they disagree, that is informative too.

\subsection*{Proxy 1: annual-report AI intensity (\var{ai\_kw}) --- main proxy for the FE model}
\begin{decisions}
\row{AI-term count per 10,000 words of the annual report, from a fixed dictionary (artificial intelligence, machine
learning, NLP, chatbot, GenAI, RPA, \ldots, and ``AI'' as an uppercase standalone token).}
{Available for almost every bank-year (387 of 400 reports), consistent across banks, and scriptable. Cost: it measures
\emph{disclosure}, not use; banks can talk about AI without using it (``AI-washing'').}
\row{``AI'' matched case-sensitively; acronyms (LLM, NLP, RPA) only in ordinary lowercase prose; exclusions for
``Air India'', law degrees (LL.M), ``LLMS'' (a loan-management system) and livestock ``artificial insemination (AI)''.}
{A match audit before any results found real false positives: e.g.\ 100 spurious ``LLM'' hits in one PNB report from a
legacy Hindi font. Fixing the dictionary before estimation keeps it a measurement decision, not a result-driven one.}
\row{Denominator = \emph{English} words only (tokens found in at least 3 private-bank reports, plus a banking-jargon
whitelist). AI terms are counted on the raw text.}
{Bilingual PSU reports in legacy Hindi fonts extract as Latin-script gibberish, inflating PSU word counts and biasing
PSU AI intensity \emph{downwards}. The whitelist stops words like CASA, GNPA, UPI being dropped. Sanity check reported:
English-only reports should drop under 15\% of tokens; above 40\% is flagged.}
\row{Robustness versions: year-standardised (z-score within FY) and $\ln(1+\text{count})$.}
{After ChatGPT (FY2023+) every bank mentions AI more; the within-year z-score asks ``which bank talks about AI more
than its peers \emph{that year}''. $\ln(1+\text{count})$ matches the source paper.}
\end{decisions}

\subsection*{Proxy 2: first deployed AI system (\var{post\_adopt}) --- main proxy for the DiD}
\begin{decisions}
\row{Adoption year = FY of the bank's first \emph{verified, deployed} AI system.}
{Deployment is closer to real adoption than talk. Alternatives kept as robustness: Alt~1 = first AI event of any type
(incl.\ partnerships); Alt~2 = first FY \var{ai\_kw} crosses and stays above the pooled median / 75th percentile.}
\row{Strict verification: every event needs a primary source (exchange filing, bank press release or annual report)
\emph{and} a dated news report. Awards, conference talks and executive interviews are not events.}
{Dates drive the DiD, so a wrong date contaminates the treatment. A single press release can be PR; two independent
sources fix the date. Cost: 88 candidates gave 54 verified events, below the 70+ target; I report the true count.}
\row{Strict rule kept even when it left banks never-treated (e.g.\ Bank of Maharashtra's 2023 ML model had a bank
release but news that did not name the model).}
{Applying the rule the same way to every bank matters more than maximising treated banks. Result: 19 of 25 banks
treated, 6 never-treated (Central, UCO, BoM, KVB, DCB; Indian Bank adopts in FY2027, after the window).}
\row{Adoption after the last panel year coded as never-treated within the window.}
{A bank first treated in FY2027 contributes no post-period data; coding it as treated would create a cohort with no
treated observations.}
\end{decisions}

\subsection*{Proxies 3--6 (robustness)}
\begin{decisions}
\row{\var{ai\_events\_cum}: cumulative verified AI events up to FY end.}
{Captures intensity of adoption, not just the first event. Cost: depends on media coverage, which is higher for large
private banks.}
\row{\var{fintech\_idx}: equal-weighted (1/6 each) RBI mobile-banking volume (min--max scaled) + 5 feature dummies
(mobile app, UPI, chatbot, voice assistant, video-KYC), following Hamdouni (2025, IJFS).}
{Measures digital infrastructure that AI builds on, from a regulator's data. Cost: feature go-live years come from the
first annual-report mention, which can be off by up to a year; it is a digital-adoption proxy, not AI per se.}
\row{\var{ai\_patents} (pre-specified) \textbf{dropped}; replaced by \var{ai\_capability}.}
{The pre-specified Google Patents search found 0 AI patent applications for all 25 banks (1 filing in total, 400 of 400
bank-years zero). A constant cannot be estimated. Kept as a descriptive fact: Indian banks buy AI, they do not patent
it. Logged as a spec change before building the replacement.}
\row{\var{ai\_capability} = sum of 4 annual-report indicators (0--4): dedicated digital/data/AI leader; AI/analytics/
innovation lab or CoE; digital/AI/data committee; in-house data-science/analytics team. As observed per report, no
carry-forward.}
{Closest available measure of ``innovative capability'' when R\&D and tech-staff data do not exist. Exclusions are
designed in: CIO/CTO/CISO (every bank has them), the RBI-mandated IT Strategy Committee, risk/audit committees,
external bodies (IBA, ministries), school CSR labs, and treasury/risk ``analytics units''. These refinements were made
after reading match snippets but before any estimate.}
\row{\var{ai\_calls}: AI terms per 10,000 words in earnings-call transcripts, FY2022 onward only. Files under
1,000 words (cover letters linking to the bank's website, scans) are excluded, so the bank-year is NA, not 0.}
{Management's unscripted attention to AI. Cost: transcripts are only filed on BSE from about 2021, so it is a
sub-period check; Axis, Kotak, IndusInd and RBL file only cover letters, leaving 21 banks. Counting cover letters as
zero AI talk would have biased these private banks downward.}
\end{decisions}
\begin{qa}
\q{Isn't counting AI words just measuring hype?}{Partly, which is why it is one of six proxies. I control for the
general rise with year fixed effects, add a within-year standardised version, and compare it with verified deployments.
The keyword measure finds no average effect, which is at least consistent with talk not translating into returns.}
\q{Why was your patent proxy dropped --- isn't that convenient?}{It was dropped because it had literally no variation:
zero AI patents for every bank-year. I logged the drop and the zero share before building the replacement, and I kept
the replacement's definition fixed before estimating it. It turned out null anyway (ROE $-0.57$ pp, $p=0.57$).}
\q{Why does Hindi text matter for a keyword count?}{PSU reports are bilingual and older ones use non-Unicode Hindi fonts
that extract as Latin gibberish. Counting those as words would make PSU banks look less AI-intensive. Using an
English-word denominator removes that bias.}
\q{Why insist on a news source if you have the bank's own filing?}{The filing proves the event; the dated news proves
\emph{when} the market learned it. For the DiD and event study, the date is the treatment.}
\end{qa}

\section{Controls and moderators}
\begin{decisions}
\row{Controls, all lagged one year: ln(total assets), capital adequacy (CRAR), loan-to-deposit ratio, liquidity
ratio, gross NPA ratio, and \var{pca} (RBI Prompt Corrective Action).}
{Lagging avoids the control being an outcome of the same year's AI use (bad control / simultaneity). PCA directly
restricts lending and expansion of weak PSU banks (2017--2021), so leaving it out would confound PSU results.}
\row{Not controls: firm age, profit growth, GDP.}
{Age is collinear with bank + year fixed effects; profit growth is mechanically tied to ROE; GDP is absorbed by year
fixed effects (used only in Phase 1). Age enters as baseline age $\times$ AI and as a DiD covariate.}
\row{PCA coded from primary sources only: exits from RBI press releases, entries from the bank's own exchange filing
quoting RBI's letter, cross-checked with the March 2018 PIB release; \var{pca}=1 if under PCA for $\ge$6 months of the
FY.}
{RBI does not issue press releases for PCA entries, so the original ``RBI press releases only'' rule could not date
entries; I changed the rule, logged it, and saved every URL. Nothing was coded from memory.}
\row{Moderators: PSU, COVID (FY2021), demonetisation (FY2017), post-ChatGPT (FY2023+), baseline age, size, baseline
Z-score (FY2011--15).}
{Tests whether AI helps more for riskier, older, or state-owned banks and in shock years. Cost: many interactions, so I
report them all and do not cherry-pick.}
\end{decisions}
\begin{qa}
\q{Why lag the controls?}{Same-year NPA or capital could themselves respond to AI (e.g.\ better credit scoring lowers
NPAs). Controlling for them would absorb part of the effect I want to measure.}
\q{Why is PCA important for this question?}{Five panel PSU banks were under PCA; they could not grow loans, which hits
ROE independently of AI. Without the control, PSU comparisons would mix AI with regulatory restrictions.}
\end{qa}

\section{Phase 2a: two-way fixed effects (main model for \var{ai\_kw})}
\[ Y_{it} = \beta\,\text{AI}_{i,t-1} + \gamma' X_{i,t-1} + \alpha_i + \delta_t + \varepsilon_{it} \]
\begin{decisions}
\row{Bank fixed effects $\alpha_i$ and year fixed effects $\delta_t$; R \var{fixest::feols}.}
{Bank FE remove time-invariant differences (ownership, franchise, management culture); year FE remove economy-wide
shocks (demonetisation, COVID, rate cycles, the general AI trend). Identification comes from a bank's AI intensity
rising faster or slower than other banks' in the same year. Cost: only within-bank variation is used.}
\row{Lag 1 is the pre-specified main; lags 0 and 2 reported.}
{AI should affect performance with a delay, and a lag reduces reverse causality (profitable banks can afford AI). Lag 0
stock return is significant (+8.19\%, wild $p=0.039$) but lag 1 is the pre-registered headline, so I do not promote
lag 0 after seeing it.}
\row{SEs clustered by bank; wild cluster bootstrap $p$-values (restricted, Webb 6-point weights, 9,999 draws).}
{With 25 clusters, conventional clustered SEs over-reject. Webb weights are designed for few clusters. The R package
for this failed with the current \var{fixest}, so I implemented the bootstrap myself on fixed-effect-demeaned data.}
\row{Effects reported per 1 SD of the proxy.}
{Makes proxies with different scales comparable (1 SD of \var{ai\_kw} = 1.47 AI terms per 10,000 words).}
\end{decisions}
\textbf{Results (lag 1, per 1 SD).} ROE +0.13 pp ($p=0.86$); stock return $-1.72\%$ ($p=0.70$); ROA +0.01 pp
($p=0.76$); Tobin's Q $-0.006$ ($p=0.29$). AI~$\times$~PSU on ROE +4.39 pp (SE 1.49, $p=0.007$, wild $p=0.007$), with the
private-bank slope $-3.20$ pp ($p=0.038$). AI~$\times$~baseline age +1.99 pp ($p=0.033$). AI~$\times$~demonetisation on
stock return $-161\%$ ($p=0.012$) is identified from a single year and I treat it as fragile.
\textbf{Post-hoc check of AI~$\times$~PSU (ROE).} Logged as post-hoc before running.
\begin{center}\small
\begin{tabular}{lccc}
\toprule
Specification & AI~$\times$~PSU & $p$ (wild $p$) & Implied PSU slope ($p$)\\
\midrule
(1) Bank + year FE (pre-specified) & +4.39 pp & 0.007 (0.006) & +1.18 pp (0.077)\\
(2) Bank + PSU-specific year FE & +1.66 pp & 0.34 (0.35) & +0.77 pp (0.24)\\
(3) Bank-specific linear trends + year FE & +3.15 pp & 0.19 (--) & +0.58 pp (0.52)\\
\bottomrule
\end{tabular}
\end{center}
Once PSU banks get their own year effects, the interaction falls by about 60\% and is no longer significant. The
pre-specified estimate stays the reported result, with this check beside it.
\begin{qa}
\q{Your hypothesis was a weaker effect for PSU banks, but the interaction is positive. How do you explain it?}{I report
it as found, and then I tested the most obvious alternative: PSU profitability recovered strongly after the AQR
clean-up from about FY2019, exactly when PSU AI disclosure rose, and year fixed effects are common to all banks. With
PSU-specific year effects the interaction drops from +4.39 to +1.66 pp ($p=0.34$), and with bank-specific trends to
+3.15 pp ($p=0.19$). So the positive sign is largely the PSU recovery trend, not robust evidence that AI helps PSU banks
more. Honest conclusion: no reliable difference between PSU and private banks.}
\q{Why the wild bootstrap?}{Inference with 25 clusters is unreliable with standard cluster-robust SEs. The wild cluster
bootstrap with Webb weights is the standard fix; here it barely changes the main $p$-values, which is reassuring.}
\q{Could profitable banks simply talk more about AI (reverse causality)?}{Yes, that is the main threat. The lag, bank
fixed effects and the DiD on deployment dates are partial answers; none fully solves it without an instrument.}
\end{qa}

\section{Phase 2b: staggered difference-in-differences (\var{post\_adopt})}
\begin{decisions}
\row{Callaway--Sant'Anna estimator (\var{did::att\_gt}), not-yet-treated controls, covariates baseline age and baseline
ln(assets); dynamic (event-time) and overall ATT.}
{Banks adopt in different years. A plain two-way FE DiD with staggered timing uses already-treated banks as controls
and can put negative weights on some effects (Goodman-Bacon). CS compares each adoption cohort only with banks not yet
treated. Cost: needs enough not-yet-treated banks in each year, which is thin here.}
\row{Sun--Abraham (\var{fixest::sunab}) as robustness; Alt~1 and Alt~2 adoption dates.}
{A second heterogeneity-robust estimator; if both agree on sign that is reassuring, but magnitudes differ here.}
\row{Post-hoc: re-estimated excluding Yes Bank, labelled post-hoc.}
{The negative ROE effect starts at event year 3, where Yes Bank's FY2020 reconstruction-year ROE ($-67.5\%$) enters.
Checking that one bank does not drive the result is reasonable; labelling it post-hoc is honest.}
\row{Post-hoc parallel-trends diagnostics (main adoption date): (a) CS on private banks only, with not-yet-treated and
never-treated private banks as controls; (b) Sun--Abraham on all banks with bank FE and PSU~$\times$~year FE.}
{The late and never-treated group is mostly PSU banks, whose ROE fell to about $-10.7\%$ (FY2018, AQR) and recovered to
about $+12.9\%$ (FY2025), so they are poor controls for early private adopters. (a) removes PSU banks as controls;
(b) keeps them but gives PSU banks their own year shocks. Cost: (a) leaves only 2--3 never- or late-treated private
banks; (b) uses the not-yet-treated PSU banks only within the PSU group.}
\end{decisions}
\textbf{Results (effect of \var{post\_adopt}=1).} Main: ROE $-23.9$ pp ($p=0.004$; Sun--Abraham $-6.9$ pp);
stock return $-37.2\%$ ($p=0.001$; Sun--Abraham $-19.1\%$). Alt~2 (keyword thresholds): not significant. Excluding Yes
Bank: ROE $-21.1$ pp ($p=0.036$), stock $-34.8\%$ ($p=0.012$). Pre-trends: three individually significant ROE
pre-period estimates (event years $-7$, $-6$, $-4$); stock-return pre-periods $-9$ to $-1$ not significant, two longer
leads ($-11$, $-10$) significant with opposite signs.

\textbf{Post-hoc parallel-trends diagnostics (main adoption).} Logged as post-hoc before running.
\begin{center}\footnotesize\setlength{\tabcolsep}{4pt}
\begin{tabular}{llccc}
\toprule
Specification & Outcome & ATT (SE, $p$) & Treated / never-late & Sig.\ leads\\
\midrule
CS, all banks (pre-specified) & ROE & $-23.9$ pp (8.4, 0.004) & 14 / 11 & 3\\
(a) CS, private banks only & ROE & $-3.27$ pp (5.82, 0.57) & 10 / 3 & 3 of 7\\
(b) Sun--Abraham, PSU~$\times$~year FE & ROE & $-2.48$ pp (1.23, 0.043) & 14 / 11 & 6 of 13 (all neg.)\\
\midrule
CS, all banks (pre-specified) & Stock return & $-37.2\%$ (11.2, 0.001) & 19 / 6 & 2 long leads\\
(a) CS, private banks only & Stock return & $-36.3\%$ (12.5, 0.004) & 11 / 2 & 0 of 7\\
(b) Sun--Abraham, PSU~$\times$~year FE & Stock return & $-15.8\%$ (5.8, 0.006) & 19 / 6 & 1 of 14\\
\bottomrule
\end{tabular}
\end{center}
Once the PSU recovery is kept out of the comparison, the ROE effect shrinks by 85--90\% and is not significant among
private banks; the remaining Sun--Abraham ROE effect sits on negative long leads (event years $-12$ to $-8$ between
$-6.3$ and $-12.2$ pp), so it is not clean either. The negative stock-return ATT survives both checks with mostly clean
leads, but the private-only version rests on two never-treated banks (DCB, KVB). The CS joint pre-test rejects
($p<0.001$) for both outcomes; a joint lead test for Sun--Abraham is not computable with 25 clusters.
\begin{qa}
\q{A $-24$ pp effect on ROE is huge. Do you believe it?}{No, not as a causal effect. The ROE pre-trends fail, the
never-treated group is six banks (mostly weak PSU banks), early adopters are large private banks, and the two
estimators differ by a factor of three. When I compare private adopters only with private banks, or give PSU banks
their own year effects, it falls to about $-3$ pp and is not significant among private banks (post-hoc). So the honest
reading is that deployment timing is correlated with bank type and the post-2016 bad-loan cycle.}
\q{Does the negative stock-return effect survive?}{In sign, yes: $-36.3\%$ among private banks and $-15.8\%$ with
PSU-specific year effects, with mostly clean short leads. But the private-only control group is two banks and the
estimators disagree on size, so I read it as a fragile association, not a cost of AI.}
\q{Why not-yet-treated rather than never-treated controls?}{The never-treated group is small and unrepresentative.
Not-yet-treated controls use more banks and compare like with like in timing.}
\q{Why did adding four new adopters not change the ROE estimate at all?}{They all adopt in FY2026, after the last
accounting year. Through FY2025 they serve as controls either way, so the ROE comparison is identical; only stock
returns (observed in FY2026) changed.}
\end{qa}

\section{Phase 3: event study}
\begin{decisions}
\row{Market model, estimation window [$-250$, $-30$] trading days, minimum 120 valid days; benchmark Nifty 500 (main),
Nifty Bank (robustness).}
{Standard short-horizon design. The 30-day gap keeps leakage out of the estimation. Nifty Bank checks whether a result is
just sector co-movement. Cost: the two benchmarks disagree for PSU banks.}
\row{CAR windows [$-1$, $+1$] (headline) and [$-5$, $+5$].}
{Short window = cleanest attribution; the longer window allows for leakage and slow diffusion.}
\row{Tests: Kolari--Pynn\"onen adjusted BMP and Corrado rank test.}
{BMP handles event-induced volatility; the Kolari--Pynn\"onen adjustment handles cross-correlation when several firms
announce on the same dates (Global Fintech Fest weeks, Microsoft India announcements in Jan 2025). Corrado is
non-parametric, robust to fat tails. Agreement across both is required before I call a result robust.}
\row{Event date = first public announcement; after 15:30 IST $\rightarrow$ next trading day.}
{The market cannot react to news released after the close.}
\row{Exclusions: earnings release within $\pm3$ trading days; crisis windows (Paytm Payments Bank action, IndusInd
derivative disclosure, Kotak and Bajaj Finance RBI embargoes, IL\&FS); pre-listing events; the firm's other events are
dropped from its estimation window.}
{Confounding news would swamp a small AI effect. Cost: fewer events (39 eligible, 30 used).}
\row{Post-hoc overlap rule: if two events of the same firm have overlapping [$-5,+5$] windows, keep only the earlier.
Paytm E070 (6 Oct 2025) is kept and E039 (7 Oct 2025) dropped.}
{The two windows share the same returns, so keeping both counts them twice and breaks the independence the tests
assume. The spec had no rule, so it was added post-hoc and logged; keeping the earlier event is mechanical (it does not
depend on which CAR is larger). The results with both events are recorded.}
\end{decisions}
\textbf{Results.} All 30 events: CAR[$-1,+1$] $+0.03\%$ (KP $p=0.84$, Corrado $p=0.58$); CAR[$-5,+5$] $-1.41\%$
($p=0.40$, $0.16$). PSU (8 events): CAR[$-5,+5$] $-2.70\%$ (KP $p=0.035$, Corrado $p=0.047$) against Nifty 500, but
$-2.83\%$ ($p=0.18$) against Nifty Bank. Private, NBFC and fintech (4 events): nothing significant at 5\%. Before the
overlap rule (31 events): all events CAR[$-1,+1$] $+0.13\%$ ($p=0.81$), CAR[$-5,+5$] $-1.28\%$ ($p=0.46$); the
conclusions do not change.
\begin{qa}
\q{Why might PSU AI announcements see negative returns?}{It is fragile --- it disappears against the bank index, so it
may reflect PSU-bank sector moves in those weeks. If real, a reading is that investors doubt PSU execution or see the
announcement as cost without near-term revenue. With 8 events I would not push it.}
\q{Why the Kolari--Pynn\"onen adjustment?}{Several firms announce at the same events (fintech fests), so their abnormal
returns are correlated; the standard BMP test then over-rejects. KP corrects for that cross-correlation.}
\end{qa}

\section{Robustness}
\begin{decisions}
\row{Winsorised outcomes (1/99); lagged profit growth added; year-standardised \var{ai\_kw}; ln(1+count); dropping
merger years; alternative proxies.}
{Each targets one specific worry: outliers, omitted profitability momentum, AI-washing, functional form, and the
mechanical effects of the 2020--21 PSU mergers.}
\row{Found and fixed a scaling bug before reporting: \var{fintech\_idx} was standardised by each bank's own SD instead
of the pooled SD.}
{Units must mean ``per 1 SD'' consistently. The corrected and earlier numbers are both recorded.}
\end{decisions}
\textbf{Results (per 1 SD at $t-1$).} All \var{ai\_kw} variants null on ROE and returns. \var{ai\_events\_cum}: ROE
$-1.85$ pp ($p=0.008$), stock $-7.06\%$ ($p=0.017$). \var{fintech\_idx}: ROE +0.61 pp ($p=0.77$), stock +11.96\%
($p=0.025$). \var{ai\_capability}: ROE $-0.57$ pp ($p=0.57$), stock +4.51\% ($p=0.10$). \var{ai\_calls}
(FY2023+ only, 21 banks): ROE $-0.52$ pp ($p=0.12$, $n=41$), stock $-15.4\%$ ($p=0.12$, $n=63$). \var{ai\_patents}: not
estimable (no variation).
\begin{qa}
\q{Your proxies disagree in sign. Which one do you trust?}{None alone. Talk (\var{ai\_kw}) and capability are null;
event counts are negative; the digital-infrastructure index is positively priced by the market. A consistent reading is
that broad digital infrastructure is valued by investors, while AI-specific signals are not yet visible in returns, and
announcement-heavy banks were not the most profitable ones over this period.}
\end{qa}

\section{Research integrity and engineering choices}
\begin{decisions}
\row{Raw downloads never edited; every transformation is a script; source URLs and download dates logged.}
{Anyone can rebuild the dataset from scratch. Hand edits are invisible and unreproducible.}
\row{Merge assertions: row counts and unique bank $\times$ year keys checked after every merge.}
{Silent duplication in merges is a classic source of wrong results.}
\row{Version control: public GitHub repo, commit after every task, milestone tags; secrets and large PDFs never
committed.}
{Full audit trail of what changed and when.}
\row{Python for data collection, PDF/text processing and panel assembly; R (\var{fixest}, \var{did}) for estimation.}
{Each language where its tooling is strongest; the estimators used are the reference implementations.}
\end{decisions}

\section{Limitations (say them before you are asked)}
\begin{itemize}
\item \textbf{Measurement.} No direct data on AI spending or use; every proxy is indirect. Keyword counts measure
disclosure; event counts depend on media coverage.
\item \textbf{Identification.} Adoption is not random: large private banks adopted early, weak PSU banks late or never,
and this coincides with the bad-loan cycle, PCA and mergers. Fixed effects remove level differences, not differential
trends.
\item \textbf{Small sample.} 25 banks, 6 never-treated, 30 usable events. Wide confidence intervals; individual
significant results can be chance given the number of tests.
\item \textbf{Event count.} 54 verified events, below the 70+ target, because of the strict two-source rule.
\item \textbf{Timing.} Most AI deployment is recent (FY2023+), leaving few post-adoption years, especially for
accounting outcomes that end in FY2025.
\end{itemize}

\section{Quick-fire answers}
\begin{description}[font=\normalfont\itshape,leftmargin=0pt,style=nextline]
\q{What is your main contribution?}{A transparent, multi-proxy measurement of AI adoption for Indian banks built
entirely from public data, with a pre-specified analysis plan, and evidence that average performance effects are not
detectable yet, with heterogeneity by ownership.}
\q{What would you do with more time?}{Mediation through
cost-to-income and NPAs (planned, deferred); a media-tone control; hand-verified deployment dates for the remaining
never-treated banks; and more post-adoption years.}
\q{What surprised you?}{That Indian banks file essentially no AI patents, and that the PSU interaction has the opposite
sign to the hypothesis.}
\q{What was the hardest data problem?}{Making annual-report text comparable across bilingual PSU reports in legacy Hindi
fonts and English-only private reports, without biasing PSU AI intensity downward.}
\q{Which result would you put on a slide?}{The null average effect of AI intensity on ROE and returns, with the
AI~$\times$~PSU interaction and its caveat, and the event-study null with the PSU benchmark sensitivity.}
\end{description}

\clearpage
\section{Question bank: everything an interviewer could ask}
Questions grouped by theme, each with a full model answer in the first person, using the project's logged numbers.
Do not memorise the wording; learn the structure (the point, the evidence, the caveat) and say it in your own words in
about a minute. Where an answer depends on facts only you know (your role, your plans), adjust it.

\newcommand{\Q}[2]{\item \textbf{#1}\par\nobreak\vspace{1pt}#2}
\newenvironment{qlist}{\begin{enumerate}[leftmargin=*,itemsep=7pt,label=\arabic*.,resume=qb]}{\end{enumerate}}

\subsection*{A. Motivation and the big picture}
\begin{qlist}
\Q{Explain your project in two minutes to a non-economist.}{Banks in India are spending a lot on artificial
intelligence --- chatbots, automated credit scoring, fraud detection. I asked a simple question: do the banks that adopt
AI actually become more profitable, or get valued more highly by the stock market? And is this weaker for government-owned
(PSU) banks than for private banks? Because nobody publishes how much AI a bank really uses, I measured it six different
ways, for example by counting how often a bank's annual report talks about AI, and by finding the date each bank first
put an AI system into actual use, confirmed by both an official document and a news report. I then followed 25 listed
banks from 2011 to 2026 and compared their profits and share returns before and after adoption, and I also checked how
share prices moved in the days around 30 AI announcements. The short answer: so far there is no reliable sign that AI
has raised profits or returns, and no robust difference between PSU and private banks.}
\Q{Why is this question important, and for whom?}{Three audiences care. Bank boards and managers are committing large,
mostly fixed technology budgets and want to know if they pay off. Investors want to know whether ``AI'' in an annual
report or a press release should change a bank's valuation. And the regulator and government care because PSU banks
hold a large share of Indian deposits; if they cannot turn technology into performance as well as private banks, the
gap between the two groups could widen. Evidence for emerging-market banks specifically is thin, so a careful answer
for India is useful even if it is a null.}
\Q{Why India, and why banks?}{India is a good laboratory for three reasons. First, the banking system is large and
digitising fast --- UPI, video-KYC, mobile banking --- so there is real variation in when banks adopt AI. Second, there is
a clean ownership split: state-owned and private banks operate under the same regulator and rules, which lets me test
the ownership question directly. Third, RBI publishes bank-wise financial tables for every bank, so profitability is
measured from one consistent source. Banks rather than all firms, because their accounts are comparable with each other
and because finance is one of the sectors where AI use cases (credit, fraud, service) are most direct.}
\Q{Why compare PSU and private banks?}{Because the economics of adoption differ. PSU banks have government ownership,
public-sector hiring and pay rules, procurement constraints and, historically, weaker balance sheets; private banks can
hire tech talent and change processes faster. So my hypothesis was that the same AI investment translates into less
performance in a PSU bank. It is also the policy-relevant comparison: if public banks lag on technology, that matters
for how the government manages them.}
\Q{What is your hypothesis, and what did you find?}{The hypothesis, fixed before estimation, was that AI adoption improves
ROE and stock returns, and that the effect is weaker for PSU banks. What I found: annual-report AI intensity has no
detectable average effect --- ROE $+0.13$ pp per standard deviation ($p=0.86$), stock returns $-1.72\%$ ($p=0.70$). The
AI~$\times$~PSU interaction on ROE was actually positive, $+4.39$ pp ($p=0.007$), the opposite of the hypothesis, but when
I give PSU banks their own year effects it falls to $+1.66$ pp ($p=0.34$), so it mostly reflects the PSU recovery after
the bad-loan clean-up. The difference-in-differences gives large negative effects, but its control group is dominated
by PSU banks with very different trends; once that is fixed the ROE effect largely disappears. AI announcements produce
no average abnormal stock return. So: no robust effect yet, and no robust PSU--private difference.}
\Q{What is new compared with existing work?}{I would not claim to be the first to study AI and bank performance. My
contribution is mainly in design and measurement for Indian banks: six proxies built entirely from public data instead
of one; adoption dates verified with a primary source plus a dated news report; a text measure corrected for a real
problem in Indian reports (bilingual PSU reports in legacy Hindi fonts that inflate word counts); heterogeneity-robust
staggered DiD rather than plain two-way fixed effects; and a pre-specified analysis plan with all 192 estimates logged,
including the nulls.}
\Q{What would change your conclusion?}{Three things. First, clean pre-trends in the difference-in-differences --- if
adopters and non-adopters moved together before adoption and diverged afterwards, I would take the DiD seriously as
causal. Second, more post-adoption years: most deployments are from FY2023 onward and accounting data end in FY2025,
so a slow effect would not show yet. Third, direct data on AI spending or usage at bank level, which would replace my
indirect proxies. If any of these pointed to a positive effect, I would update.}
\Q{What is the one-line takeaway?}{So far, AI talk and AI announcements have not shown up in Indian banks' profitability
or stock returns in any robust way, and the apparent PSU advantage is the post-clean-up recovery rather than AI; the
only positive association is that investors seem to reward broad digital infrastructure.}
\end{qlist}

\subsection*{B. Banking and finance domain knowledge}
\begin{qlist}
\Q{Define ROE and ROA. Why use average equity?}{ROE is profit after tax divided by shareholders' equity, times 100; ROA
is profit after tax divided by total assets, times 100. I use the average of opening and closing equity (and assets)
because profit is earned over the whole year. This matters a lot for Indian PSU banks, which received large government
capital infusions; dividing a year's profit by the post-infusion closing equity would understate the return.}
\Q{What is Tobin's Q and how did you measure it for banks?}{Tobin's Q is the market value of a firm relative to the
book value of its assets --- above 1 means the market values the bank above its balance sheet, for example for its
franchise or growth. For banks I use (market capitalisation + total liabilities) / total assets. Bank liabilities are
mostly deposits and borrowings whose book value is close to market value, so the main market component is equity. It
is noisy, which is why it is a secondary outcome.}
\Q{What are gross NPAs, CRAR and the loan-to-deposit ratio?}{The gross NPA ratio is non-performing loans --- loans more
than 90 days overdue --- as a percentage of gross advances; it measures asset quality. CRAR, the capital to
risk-weighted assets ratio, is regulatory capital as a percentage of risk-weighted assets; it measures how much loss
the bank can absorb. The loan-to-deposit ratio is advances divided by deposits; it measures how much of its funding the
bank lends out. All three, plus size, liquidity and PCA status, enter my models as one-year-lagged controls.}
\Q{What was the Asset Quality Review (AQR) and why does it matter for you?}{The AQR was RBI's inspection exercise in
2015--16 that forced banks to recognise loans they had been ``evergreening'' as non-performing. Reported bad loans
jumped --- PSU banks' gross NPA ratio rose to roughly 14--15\% around March 2018 --- and provisioning wiped out profits; PSU
banks had negative ROE for several years and then recovered from about FY2019. For me it is the single biggest
confounder: PSU profitability fell and recovered for reasons unrelated to AI, at the same time as AI disclosure rose. That
is why I checked the PSU interaction with PSU-specific year effects, and the DiD with private-only controls.}
\Q{What is Prompt Corrective Action (PCA)?}{PCA is RBI's framework for weak banks. When capital, asset-quality or
profitability thresholds are breached, RBI restricts lending, branch expansion, dividends and management pay until the
bank recovers. Five banks in my panel --- Bank of India, Bank of Maharashtra, UCO, Indian Overseas Bank and Central Bank
--- were under PCA at various points between about 2017 and 2021. Because PCA mechanically restricts growth and
profitability, I coded it from primary sources and include it as a lagged control.}
\Q{How did demonetisation and COVID affect banks?}{Demonetisation in November 2016 (FY2017) withdrew most high-value notes;
banks received a surge of deposits and digital payments got a big push. COVID in FY2021 brought a loan moratorium,
restructuring, higher provisioning and very volatile share prices. In my panel models both are absorbed by year fixed
effects, since they hit all banks in the same year; I also interact them with AI as pre-specified moderators, and use
them as dummies in the sector-level Phase 1 regressions.}
\Q{What happened to Yes Bank in 2020?}{Yes Bank had grown aggressively, with heavy exposure to stressed corporate
borrowers. In March 2020 RBI imposed a moratorium and superseded the board, and a reconstruction scheme brought in SBI
and other banks as investors. Its FY2020 ROE was about $-67.5\%$. Because it is an early adopter in my data, it sits in
the post-adoption period of the DiD and can drag the average down, so I re-ran the DiD without it, labelled as post-hoc.
The estimate barely moved ($-21.1$ pp instead of $-23.9$ pp), so Yes Bank is not what drives it.}
\Q{Why did PSU banks merge in 2019--2021, and how do mergers affect your data?}{The government consolidated PSU banks to
create fewer, stronger ones: SBI absorbed its associates in FY2018, Bank of Baroda merged with Dena and Vijaya in
FY2020, and in FY2021 PNB, Canara, Union Bank and Indian Bank each absorbed one or two banks. In the data a merger makes
assets, ratios and sometimes ROE jump mechanically. I handle this with pre-specified merger-year dummies and a robustness
check that drops merger years; the \var{ai\_kw} results do not change (ROE $+0.04$ pp, $p=0.95$, without merger years).}
\Q{What is UPI and why include it in a digital index?}{UPI, the Unified Payments Interface run by NPCI since 2016, is
India's real-time mobile payment system. Going live on UPI is a basic digital capability for a bank, and AI applications
like fraud detection and personalised offers sit on top of that digital layer. So my fintech index includes UPI, along
with a mobile app, chatbot, voice assistant, video-KYC and RBI mobile-banking volumes, as a measure of digital
infrastructure --- I describe it as a digital-adoption proxy, not as AI itself.}
\Q{How would AI actually raise bank profits?}{Through two main channels. The efficiency channel: automation, chatbots and
document processing reduce operating costs, which would show up in a lower cost-to-income ratio. The credit channel:
better credit scoring and early-warning models reduce bad loans and provisioning, which would show up in lower NPAs.
There are also revenue channels --- cross-selling, faster loan decisions, fraud prevention. Cost-to-income and NPAs were
pre-specified as mediators; testing them is deferred until after this round.}
\Q{Why might AI \emph{not} raise profits in the short run?}{Several reasons. The costs are upfront and fixed --- data
infrastructure, vendors, staff --- while benefits arrive slowly as processes change. Many ``deployments'' are pilots or
customer-facing chatbots with little effect on the core business. If all banks adopt similar tools, competition passes
the savings to customers rather than shareholders. And there is the productivity-paradox pattern seen with earlier
technologies: gains appear years later. My panel has few post-adoption years, so a delayed effect would not be visible.}
\end{qlist}

\subsection*{C. Data and sample}
\begin{qlist}
\Q{Where does every variable come from?}{Financials and controls --- profit, equity, assets, NPAs, CRAR, deposits, advances
--- come from RBI's \emph{Statistical Tables Relating to Banks in India} on DBIE. Stock prices are adjusted closing prices
from Yahoo Finance (\texttt{.NS} tickers), and the Nifty 500 and Nifty Bank indices for the event study. AI events,
earnings-call transcripts and earnings dates come from BSE/NSE corporate filings and bank press releases, with a dated
news report for each event. The keyword and capability measures come from annual-report PDFs. PCA dates come from RBI
press releases and banks' filings; mobile-banking volumes from RBI's monthly bank-wise releases; GDP from MOSPI/RBI. Every
download's URL and date is logged.}
\Q{Why only free data? Wouldn't Prowess be better?}{Prowess would have saved time, for example a ready market-cap
series. I chose public data deliberately so that anyone can rebuild the whole dataset and check my work, which matters
for a project whose credibility rests on transparency. The cost is that I had to construct some variables myself ---
market capitalisation, for example --- and that some things, like tech headcount or IT spending, are simply not
available.}
\Q{How many banks, years and observations?}{25 listed banks --- 12 PSU and 13 private --- from FY2011 to FY2026, giving
392 bank-years in an unbalanced panel. Accounting variables run to FY2025 and stock returns to FY2026. For the keyword
measure I have 387 of 400 possible annual reports. The event study uses 14 firms --- banks, NBFCs and fintechs --- and
30 announcements after all exclusions.}
\Q{Why does the accounting data stop at FY2025?}{Because RBI had not yet published the FY2026 edition of its bank-wise
tables when I built the data. I could have filled FY2026 from banks' own annual reports, but that would mix two sources
with slightly different definitions in the final year, exactly where many adoptions happen. I preferred one consistent
source and logged the decision as a spec change; stock returns, which do not depend on RBI, run to FY2026.}
\Q{How did you handle missing data?}{A strict rule: never invent or impute. If a value is missing or looks wrong, it stays
missing and is logged in the data notes with the reason. For example, 39 bank-years where RBI's paid-up capital includes
government preference capital would overstate shares outstanding, so their market cap is missing; a suspicious
Karur Vysya FY2012 return and unreliable Yahoo prices before May 2011 are also set to missing. A wrong number is worse
than a missing one because it biases results silently.}
\Q{Give an example of a data error you caught.}{The clearest one is the Hindi-font problem. Older PSU annual reports are
bilingual, and the Hindi is set in legacy non-Unicode fonts that extract from the PDF as Latin-letter gibberish. That
inflated PSU word counts and even produced fake keyword hits --- one PNB report had about 100 spurious ``LLM'' matches.
I caught it by auditing matches before estimating, and fixed it with an English-word denominator and context rules.
Others: preference capital inflating share counts, earnings-call ``transcripts'' that were really one-page cover letters,
and a scaling bug where the fintech index was standardised by each bank's own SD instead of the pooled SD.}
\Q{How did you build market capitalisation?}{No free historical market-cap series exists, so I built it: shares
outstanding = paid-up equity capital from RBI divided by face value, adjusted for face-value splits using split
records, times the fiscal-year-end share price. I validated it against Yahoo's current market cap for five banks and
required agreement within 5\%. Where paid-up capital includes preference capital, the share count would be wrong, so I
set those bank-years to missing rather than guess.}
\Q{How did you make sure merges were correct?}{After every merge the scripts assert the expected row count and that
bank~$\times$~year keys are unique, and print diagnostics --- how many rows matched, which did not. Silent duplication in
a merge is one of the most common ways to get wrong regression results, so the scripts stop rather than continue if an
assertion fails.}
\Q{Why is IDFC First only from FY2019?}{IDFC First Bank was formed in December 2018 by the merger of IDFC Bank and
Capital First, a non-bank lender. Before FY2019 it is a different entity, so stitching the earlier series on would
compare unlike things. It enters the unbalanced panel from FY2019, as the spec says.}
\Q{How did you handle the fiscal year?}{Everything uses the Indian fiscal year, April to March: FY2017 means 1 April 2016
to 31 March 2017. Banks report on this basis and RBI tables use it, so I also measure annual stock returns from 1 April
to 31 March and assign events to fiscal years the same way. That keeps accounting, market and adoption timing aligned.}
\end{qlist}

\subsection*{D. Measuring AI adoption}
\begin{qlist}
\Q{How do you define AI adoption?}{There is no administrative data on bank AI use in India, so I use six proxies, each
capturing a different facet. Two are main: \var{ai\_kw}, AI terms per 10,000 English words in the annual report
(disclosure), and \var{post\_adopt}, which switches on from the year of the bank's first verified, deployed AI system
(deployment). Four are robustness checks: cumulative verified AI events, a fintech/digital index, an AI-capability index
from annual reports, and AI intensity in earnings-call transcripts. The pre-specified patent proxy had to be dropped
because no bank had any AI patents.}
\Q{Isn't a keyword count just measuring hype or AI-washing?}{Partly, yes --- it measures disclosure, not use, and I say
so. I deal with it in three ways. Year fixed effects remove the general rise in AI talk that affects every bank, for
example after ChatGPT. A year-standardised version asks whether a bank talks about AI more than its peers in the same
year. And I do not rely on it alone: verified deployments and other proxies sit beside it. The keyword measure shows no
effect anyway, which fits the idea that talk alone does not move performance.}
\Q{Show me your dictionary. Why those terms?}{The dictionary was fixed in the spec before any estimation: artificial
intelligence, machine learning, deep learning, neural network, natural language processing and NLP, chatbot,
conversational AI, generative AI and GenAI, large language model and LLM, computer vision, robotic process automation
and RPA, predictive analytics, AI/ML, and ``AI'' as a standalone uppercase word. These are the standard terms banks
use when they describe AI projects. Fixing it in advance means I could not tune the list towards a result.}
\Q{How did you deal with false positives?}{By auditing a sample of matches in context before estimating anything. That
showed specific problems, which I fixed with rules: ``AI'' is matched only as an uppercase standalone word, and ``Air
India'' is excluded; acronyms like LLM, NLP and RPA count only inside ordinary lowercase prose, which removes Hindi-font
gibberish; ``LL.M'' (a law degree in director profiles), ``LLMS'' (a loan lifecycle management system) and
``artificial insemination (AI)'' in agricultural-lending sections are excluded. All of this was decided before results,
so it is a measurement fix, not a result-driven one.}
\Q{Why normalise per 10,000 words, and why English words only?}{Annual reports range from short to several hundred
pages, so raw counts would mostly measure report length; a rate per 10,000 words fixes that. The denominator is English
words only because bilingual PSU reports contain Hindi in legacy fonts that extracts as Latin gibberish. Counting those
tokens as words would inflate PSU denominators and make PSU banks look less AI-intensive than they are. I count a token
as English if it appears in at least three private-bank reports, plus a banking-jargon whitelist (CASA, GNPA, UPI, KYC
and so on), and I check that English-only reports lose less than 15\% of tokens to the filter.}
\Q{Is RPA really AI?}{Strictly, no --- robotic process automation is mostly rule-based scripting of repetitive tasks. It is
in the dictionary because the spec included it and because banks usually discuss it as part of their ``AI and
automation'' programme, so it tracks the same strategic intent. But I did not let an RPA-only rollout count as the first
deployed AI system for dating adoption, so it does not drive the DiD.}
\Q{How did you verify an AI event?}{Each candidate event needed two things: a primary source --- an exchange filing, the
bank's own press release or its annual report --- and an independent, dated news report. It also had to be a launch, a
partnership or vendor deal, an AI investment, or an internal deployment; awards, conference talks and executive
interviews do not count. From 88 candidates, 54 met the rule. That is below the 70+ I originally targeted, and I report
the true number rather than loosen the rule.}
\Q{Why ``first deployed'' rather than first announced?}{Because the research question is about adopting AI, and a system in
actual use is much closer to adoption than an MoU or a statement of intent; announcements are often never followed
through. But I do not throw the announcement information away: the Alt~1 definition dates adoption at the first AI
event of any type, and it gives very similar results (ROE $-25.2$ pp versus $-23.9$ pp), so the choice does not drive
anything.}
\Q{Why did you drop patents?}{The spec pre-specified an AI-patent proxy, and I ran the search exactly as planned on Google
Patents' public BigQuery data for all 25 banks. It returned zero AI patent applications for every bank-year --- one
patent filing of any kind in total. A variable that is always zero cannot be estimated. I logged that before building a
replacement, and I keep it as a descriptive finding: Indian banks buy AI from vendors, they do not patent it.}
\Q{How is the AI-capability index built, and how did you avoid false matches?}{It is the sum of four yes/no indicators
read from each annual report: a dedicated digital, data or AI leader (such as a Chief Digital or Data Officer); an AI,
analytics or innovation lab or centre of excellence; a board or management committee specifically for digital, AI or
data; and an in-house data-science or analytics team. So it runs from 0 to 4. To avoid false matches I excluded roles
every bank has (CIO, CTO, CISO), the RBI-mandated IT Strategy Committee, risk and audit committees, external bodies
like IBA, school CSR labs, and treasury ``analytics units''. Those refinements were made after reading match snippets
but before any estimate.}
\Q{Why equal weights in the fintech index?}{I followed Hamdouni (2025) in the \emph{International Journal of Financial
Studies}, which builds a comparable index with equal weights. Any other weighting would be my own arbitrary choice, and
with only six items there is no strong case for one. The cost is that a scaled transaction volume and a 0/1 feature
dummy count equally; I accept that and say it.}
\Q{Why use earnings-call transcripts, and what went wrong?}{Earnings calls are unscripted, so AI mentions there reflect
management attention more than a polished annual report does. Two limits: transcripts are filed on BSE only from
around FY2022, so this is a sub-period check, and four banks --- Axis, Kotak, IndusInd and RBL --- file only a one-page
cover letter linking to the transcript on their website. At first those were being read as transcripts with zero AI
words, which would bias these private banks down. I excluded files under 1,000 words, so those bank-years are missing
rather than zero, leaving 21 banks.}
\Q{Which proxy do you trust most?}{It depends on the job. For timing --- when did a bank adopt --- I trust verified
deployment most, because it is dated from two sources. For coverage and a continuous measure, the annual-report keyword
intensity, which exists for almost every bank-year. But I would not rest a conclusion on any single proxy; that is the
point of having six.}
\end{qlist}

\subsection*{E. Econometrics: fixed effects}
\begin{qlist}
\Q{Write down your main model and interpret $\beta$.}{$Y_{it}=\beta\,\text{AI}_{i,t-1}+\gamma'X_{i,t-1}+\alpha_i+
\delta_t+\varepsilon_{it}$, where $Y$ is ROE (or stock return, ROA, Tobin's Q), AI is \var{ai\_kw} scaled per standard
deviation and lagged one year, $X$ are lagged controls, $\alpha_i$ bank effects and $\delta_t$ year effects. $\beta$ is
the change in ROE, in percentage points, when a bank's AI intensity is one SD (1.47 AI terms per 10,000 words) higher
than its own average, relative to how other banks moved that year. My estimate is $+0.13$ pp with SE 0.73, so I cannot
distinguish it from zero.}
\Q{What do bank and year fixed effects control for, and what do they not?}{Bank effects remove everything about a bank that
does not change over time --- ownership, franchise, region, culture. Year effects remove shocks common to all banks in a
year --- demonetisation, COVID, interest-rate cycles, the general rise in AI talk. What they do not remove are
bank-specific changes over time that are correlated with AI: for example, PSU banks recovering after the AQR exactly
when they started talking about AI. That is why I added PSU-specific year effects and bank-specific trends as checks.}
\Q{Why cluster by bank? Why not by year or two-way?}{The errors of the same bank are correlated across years --- a bank with
an unusually good year tends to have another --- so standard errors must allow for that, which clustering by bank does.
Clustering by year would need many years; I have about 15, which is too few clusters to be reliable, and two-way
clustering would inherit that problem. With 25 bank clusters, even bank clustering is on the low side, so I add wild
cluster bootstrap $p$-values.}
\Q{What is a wild cluster bootstrap and why Webb weights?}{With few clusters, conventional cluster-robust standard errors
are too small and reject too often. The wild cluster bootstrap fits the model with the null hypothesis imposed, then
creates many artificial samples by multiplying each bank's residuals by a random weight --- the same weight for all of
that bank's years --- and re-estimates the $t$-statistic each time. The $p$-value is the share of artificial
$t$-statistics larger than the real one. Webb weights take six values instead of Rademacher's two ($\pm1$), which gives
more distinct bootstrap draws when clusters are few. I used 9,999 draws.}
\Q{Why did you write your own bootstrap?}{The standard R package for the wild cluster bootstrap failed with the current
version of \var{fixest}, so I implemented the restricted wild cluster bootstrap myself: demean the data for the fixed
effects, impose the null, apply Webb weights by bank, and compute cluster-robust $t$-statistics. As a check, it gives
$p$-values close to the conventional ones where they should be close (ROE: 0.86 vs 0.86; AI~$\times$~PSU: 0.007 vs 0.007).}
\Q{Why lag AI by one year? Why report lags 0 and 2?}{AI projects take time to affect costs or loan quality, so a lag
is economically natural, and it reduces reverse causality, because this year's profit cannot cause last year's AI
disclosure. Lag 1 was pre-specified as the main result. I report lags 0 and 2 for transparency: at lag 0 the stock
return effect is $+8.19\%$ ($p=0.006$, wild $p=0.039$), which could be the market reacting to same-year disclosure,
and at lag 2 ROE is $-2.85$ pp ($p=0.10$). Neither changes the main conclusion, and I do not switch to whichever lag is
significant.}
\Q{Interpret the AI~$\times$~PSU coefficient. What is the PSU slope?}{With the interaction, the AI coefficient alone is
the slope for private banks: $-3.20$ pp of ROE per SD of AI intensity ($p=0.038$). The interaction, $+4.39$ pp, is how
much larger the slope is for PSU banks. So the PSU slope is $-3.20+4.39=+1.18$ pp per SD (SE 0.64, $p=0.077$): slightly
positive for PSU banks, negative for private banks. That is the opposite of my hypothesis --- which is why I tested
whether it is a trend artefact.}
\Q{Is the PSU result driven by the post-AQR recovery?}{Largely, yes. PSU profitability collapsed after the AQR and
recovered from about FY2019, exactly when PSU banks began writing more about AI, and ordinary year effects are common to
both groups, so the recovery can load onto the interaction. When I let PSU banks have their own year effects, the
interaction falls from $+4.39$ to $+1.66$ pp ($p=0.34$); with bank-specific linear trends it is $+3.15$ pp ($p=0.19$).
This check is post-hoc and labelled as such, and the pre-specified estimate is still reported, but my honest reading is
that there is no reliable PSU--private difference.}
\Q{How many tests did you run, and how do you handle multiple testing?}{The spec log has 192 estimated specifications ---
proxies, outcomes, lags, interactions, robustness checks. At a 5\% level, about ten would be significant by pure chance.
I handle it by pre-specifying the main estimates (lag 1, ROE and stock return, \var{ai\_kw} and \var{post\_adopt}),
reporting every estimate including nulls, and treating isolated significant results as fragile --- for example the
AI~$\times$~demonetisation interaction, which is identified from a single year. I would rather show the reader all the
tests than adjust a hand-picked subset.}
\Q{What is reverse causality here and how do you address it?}{Profitable banks can afford more AI and may also write
more about it, so a correlation between AI and ROE could run from profits to AI. I address it partly: AI is lagged, so
current profits cannot cause it; bank fixed effects remove the fact that some banks are always more profitable; and the
DiD uses dated first deployments. None of these fully solves it, because a bank expecting good years may invest first.
I say that plainly.}
\Q{Why not use an instrument?}{Because I could not find a credible one. An instrument has to shift AI adoption but affect
profits only through AI. Candidates like IT budgets, management quality or vendor exposure all affect profits directly;
regional digital infrastructure affects deposit growth directly. A weak or invalid instrument would give a more
precise-looking but less trustworthy answer, so I prefer to be transparent about the limits of fixed effects and DiD.}
\Q{Why not include firm age as a control?}{Because with bank and year fixed effects, age is perfectly collinear: age =
year minus founding year, and both pieces are already absorbed by the fixed effects. Instead, age enters where it can be
identified --- baseline age (age in FY2011) interacted with AI, and as a covariate in the DiD.}
\end{qlist}

\subsection*{F. Econometrics: difference-in-differences}
\begin{qlist}
\Q{What is the problem with two-way FE DiD under staggered adoption?}{When banks adopt in different years, the ordinary
two-way fixed-effects DiD is a weighted average of many two-by-two comparisons, and some of those use already-treated
banks as the control group for later adopters --- the ``forbidden comparisons'' in the Goodman-Bacon decomposition. If
effects grow over time, those comparisons can get negative weights, so the estimate can even have the wrong sign. That
is why I use estimators built for staggered timing.}
\Q{Explain Callaway--Sant'Anna in plain words.}{Group banks by the year they first adopted --- a cohort. For each cohort and
each year, compare how that cohort's outcome changed since just before adoption with how banks that had not yet adopted
changed over the same period. That gives one clean effect per cohort-year, never using already-treated banks as
controls. Then average those effects --- overall, or by years since adoption for an event-study plot. I also condition on
baseline age and baseline size using the doubly-robust version.}
\Q{Never-treated vs not-yet-treated controls?}{Never-treated controls are banks that never adopt in the sample;
not-yet-treated also include future adopters before they adopt. I pre-specified not-yet-treated because my never-treated
group is only six banks, mostly weaker PSU banks, which is small and unrepresentative; using future adopters adds more
comparable banks. Both choices still lean heavily on PSU banks, which is the problem the post-hoc diagnostics address.}
\Q{What is the parallel-trends assumption, and does it hold?}{It says that, without adoption, adopters and their control
banks would have followed the same trend in the outcome. It cannot be tested directly, but pre-adoption estimates should
be close to zero. For ROE they are not: three pre-period estimates are individually significant, and the joint
pre-test rejects. The reason is visible in the data: early adopters are large private banks, while the late and
never-treated controls are mostly PSU banks whose ROE fell to about $-10.7\%$ in FY2018 and recovered to about $+12.9\%$
by FY2025. So for ROE, parallel trends is not credible in the main specification.}
\Q{What did the post-hoc DiD diagnostics show?}{I ran two checks, labelled post-hoc, with the main adoption definition.
First, Callaway--Sant'Anna on private banks only: ROE $-3.27$ pp ($p=0.57$), stock return $-36.3\%$ ($p=0.004$),
with 10--11 treated and only 2--3 control banks. Second, Sun--Abraham on all banks with PSU-specific year effects: ROE
$-2.48$ pp ($p=0.043$), but on negative long leads, and stock return $-15.8\%$ ($p=0.006$). So the large $-23.9$ pp ROE
effect was mostly the PSU control group. The negative stock-return effect survives in sign, but rests on very few
controls, so I treat it as a fragile association.}
\Q{Why do CS and Sun--Abraham differ so much ($-23.9$ vs $-6.9$ pp)?}{They weight and compare differently.
Callaway--Sant'Anna uses not-yet-treated banks as controls, with covariate adjustment, and its overall ATT gives weight
to long horizons where only the earliest cohorts --- including Yes Bank's FY2020 collapse --- are observed. Sun--Abraham
compares cohorts with the never-treated group, includes time-varying controls, and weights cohort effects by cohort
size. With few cohorts, few control banks and one extreme bank-year, the estimates are sensitive to those choices. I
report both; the disagreement is itself a warning not to read the magnitude causally.}
\Q{Why did you exclude Yes Bank, and isn't that cherry-picking?}{After looking at the event-time plot I saw the negative ROE
effect start at event year 3, exactly where Yes Bank's $-67.5\%$ FY2020 ROE enters. Checking whether one bank drives a
result is a reasonable robustness question, but since I asked it after seeing results, it is labelled post-hoc and the
pre-specified full-sample estimate stays the headline. The answer was that it does not: ROE $-21.1$ pp ($p=0.036$)
without Yes Bank, versus $-23.9$ pp with it.}
\Q{What are the covariates in the DiD and why?}{Baseline age and baseline log total assets, both measured at the start of
the sample. Bigger and older banks plausibly adopt AI at different times and also follow different profit trends, so
assuming parallel trends \emph{conditional} on these is more plausible than unconditionally. They are measured at
baseline so they cannot be affected by adoption itself.}
\Q{How would you improve the DiD?}{First, more post-adoption years, which only time solves. Second, verified deployment
dates for the remaining never-treated banks, to get more and better controls. Third, a formal sensitivity analysis such
as Rambachan and Roth's HonestDiD, which shows how large a violation of parallel trends would overturn the result.
And with more private banks, the private-only design could become the main one.}
\end{qlist}

\subsection*{G. Event study}
\begin{qlist}
\Q{Walk me through an event study.}{For each announcement, I first estimate what the stock's ``normal'' return should
be: a market model regressing the stock's daily return on the Nifty 500 over trading days $-250$ to $-30$ before the
event, requiring at least 120 valid days. Then, around the event, the abnormal return each day is the actual return minus
the model's prediction. I add these up over a window --- $[-1,+1]$ as the headline, $[-5,+5]$ as a wider check --- to get
the cumulative abnormal return, or CAR, and test whether the average CAR across events differs from zero.}
\Q{Why the market model and not CAPM or Fama--French?}{For windows of a few days the choice of normal-return model makes
very little difference, because expected returns over three days are tiny relative to the abnormal returns of interest;
the market model is the standard choice in short-window studies. CAPM imposes restrictions the market model does not
need, and multi-factor Indian data are less standardised and less freely available. So the market model is both
standard and adequate here.}
\Q{Why two benchmarks?}{The Nifty 500 represents the whole market, so it is the main benchmark. The Nifty Bank index
represents the banking sector; if a result appears against the market but not against the sector index, the ``abnormal''
return may just be the banking sector moving that week. That is exactly what happens with my PSU result, which is why I
call it fragile.}
\Q{Why BMP, Kolari--Pynn\"onen and Corrado tests?}{Each handles a different problem. The BMP test standardises abnormal
returns, so it is robust to the higher volatility that often comes with events. The Kolari--Pynn\"onen adjustment
corrects BMP for correlation across events, which matters because several firms announced on the same dates --- Global
Fintech Fest weeks, the January 2025 Microsoft announcements --- and correlated returns make the plain test over-reject.
The Corrado rank test is non-parametric, so it does not assume normal returns, and stock returns have fat tails. I
call a result robust only if both tests agree.}
\Q{How did you choose the event date?}{The event date is the first public announcement, from the exchange filing or press
release, cross-checked with the dated news report. If it was released after the market closed at 15:30 IST, it moves to
the next trading day, because the market could not react before then. Filing timestamps were used for this where
available.}
\Q{Why exclude events near earnings releases and crisis windows?}{Because the AI effect, if any, is small, and an earnings
release or a crisis moves prices far more. If an AI announcement falls within three trading days of results, I cannot
separate the two. Likewise for crisis windows like RBI's action on Paytm Payments Bank, the embargoes on Kotak and Bajaj
Finance, IndusInd's derivative-accounting disclosure and the IL\&FS crisis. Excluding them costs events but keeps the
CARs interpretable.}
\Q{How many events did you use and why so few?}{54 verified events, of which 39 were eligible for CARs --- the others were
for unlisted periods, such as before Paytm's or SBI Cards' listing, or for panel-only banks. After excluding events near
earnings and in crisis windows, and applying the same-firm overlap rule, 30 remain. It is few because I insisted on a
primary source plus a dated news report for every event; a looser rule would give more events but noisier dates.}
\Q{What about overlapping events for the same firm?}{Two Paytm events were one trading day apart (E070 on 6 October 2025
and E039 on 7 October 2025), so their $[-5,+5]$ windows shared the same returns and would be counted twice. The spec had
no rule for this, so I added one, logged as post-hoc: when same-firm windows overlap, keep only the earlier event. It is
mechanical --- it does not depend on which CAR is larger. E039 was dropped, and the results barely changed: all-event
CAR$[-1,+1]$ moved from $+0.13\%$ to $+0.03\%$.}
\Q{Why could markets react negatively to AI news?}{Several reasons. Investors may see an announcement as cost without
near-term revenue. For PSU banks they may doubt execution. The news may already be priced in because it leaked or was
expected, so the announcement day shows nothing or a ``sell the news'' move. And some announcements are PR for a weak
quarter. My average CARs are close to zero, so I do not have a strong negative effect to explain except the fragile PSU
result.}
\Q{Is the PSU CAR result robust?}{No. For the eight PSU events, CAR$[-5,+5]$ is $-2.70\%$, significant against the
Nifty 500 under both tests (Kolari--Pynn\"onen $p=0.035$, Corrado $p=0.047$). But against the Nifty Bank index it is
$-2.83\%$ with $p=0.18$ and $0.09$. So it may just reflect PSU-bank sector moves in those weeks, and with eight events I
would not build an argument on it.}
\end{qlist}

\subsection*{H. Interpreting the results}
\begin{qlist}
\Q{Your main proxy is null. Is the project a failure?}{No. The question was whether AI adoption improves performance,
and ``not detectably, so far'' is an answer, especially given how much is being spent. The null is reported with
confidence intervals, not just a $p$-value, and it holds across winsorising, alternative controls, year-standardised and
log versions of the keyword measure, and dropping merger years. The project also produced a reusable dataset and
measurement method. A pre-registered null is more useful than a fragile positive result.}
\Q{Why do proxies give different signs?}{Because they measure different things. Keyword intensity measures talk; event
counts measure announcement activity, which is higher for large private banks that had weak relative returns in parts
of this period; the fintech index measures digital infrastructure; capability measures organisation. It would be more
worrying if six imperfect proxies of different concepts all agreed perfectly. The pattern I see is: AI-specific signals
are null or negative, and broad digital infrastructure is positively associated with stock returns.}
\Q{Is the \var{fintech\_idx} stock-return effect causal?}{No. It is a robustness association: one SD higher fintech index
goes with $+11.96\%$ higher annual stock return ($p=0.025$), with no effect on ROE ($+0.61$ pp, $p=0.77$). It could
reflect overall digital leadership, management quality or growth expectations rather than the features themselves, and
it is one significant result among many tests. I would describe it as suggestive.}
\Q{Why is \var{ai\_events\_cum} negative?}{One more cumulative AI event goes with lower ROE ($-1.85$ pp per SD, $p=0.008$)
and lower returns ($-7.06\%$, $p=0.017$). I read it as an association, not a cost of AI: event counts accumulate fastest
for large, heavily covered banks and in waves of PSU announcements in years when returns were weak, and the count only
rises over time, so it partly tracks bank-specific trends. It is a robustness proxy, and the main proxies do not show a
negative effect.}
\Q{What is the economic magnitude of your PSU interaction?}{$+4.39$ pp of ROE per one SD of AI intensity, when Indian
banks' ROE is typically around 10--15\% in normal years --- so it would be large if causal. That size is itself a warning
sign: it is implausible that AI talk alone moves ROE that much, and indeed it shrinks to $+1.66$ pp and becomes
insignificant once PSU-specific year effects absorb the post-AQR recovery.}
\Q{Statistical vs economic significance: give an example from your results.}{The AI~$\times$~demonetisation interaction
on stock returns is statistically significant ($p=0.012$) but implausibly large, $-161\%$, and identified from a single
fiscal year, so it means very little. The opposite case is the fintech index: $+12\%$ per SD of stock return is
economically large but only modestly significant. The lesson is to look at magnitude, identifying variation and
robustness, not just the $p$-value.}
\Q{How do your results compare with what you expected?}{I expected a positive effect that is weaker for PSU banks. I
found no robust average effect, and a PSU interaction with the opposite sign that disappears when PSU-specific trends
are controlled. The DiD looked strongly negative but turned out to be a control-group problem. So the data did not
support the hypothesis, and I reported that rather than search for a specification that would.}
\end{qlist}

\subsection*{I. Robustness, integrity and reproducibility}
\begin{qlist}
\Q{What robustness checks did you run and why each?}{Each targets one worry. Winsorising outcomes at 1/99 percent --- are
outliers driving it? Adding lagged profit growth --- is it omitted profitability momentum? Year-standardised \var{ai\_kw}
--- is it the post-ChatGPT surge in AI talk? $\ln(1+\text{count})$ --- is it functional form? Dropping merger years --- is
it the mechanical effect of the PSU mergers? And the four alternative proxies --- is it the measure? The \var{ai\_kw}
result is null in every version (ROE between $-0.80$ and $+0.27$ pp, all $p\geq0.40$).}
\Q{How did you avoid specification searching?}{By writing the analysis spec before estimating anything --- outcomes,
proxies, models, lags, controls and robustness checks. Any change had to be logged in a spec-changes file with the date
and reason before re-running, and analyses added after seeing results are labelled post-hoc everywhere. Every estimate,
including nulls, goes into the spec log, so a reader can see everything I ran, not just what I chose to show.}
\Q{Give an example of a decision you changed and why.}{Three examples. PCA dating: the spec said RBI press releases only,
but RBI issues releases only for exits, not entries, so I added banks' own exchange filings quoting RBI's letter. The
patent proxy: it had no variation at all, so I dropped it and built the capability index, logging the drop first. And
the transcripts: I found cover letters being counted as transcripts with zero AI words, so I excluded files under 1,000
words. Each was logged with the reason, and none was chosen because of how results looked.}
\Q{Could someone else reproduce your results?}{Yes. All data are public, every transformation is a script --- raw files
are never edited by hand --- each download's source URL and date is recorded, and the full history is in a public Git
repository with a commit per task and milestone tags. Annual-report PDFs and raw price dumps are not in the repository
because of size, but the scripts and URLs to re-download them are.}
\Q{What did you do when a tool or package failed?}{I worked around it transparently. When the wild-bootstrap package failed
with the current \var{fixest}, I implemented the bootstrap myself and checked it against conventional $p$-values. When
service-account keys for BigQuery were blocked, I ran the patent query in the BigQuery web console with the same SQL.
In both cases the method stayed the pre-specified one; only the tooling changed.}
\Q{How did you validate your text-mining?}{I audited matches in context before estimating: scripts write snippet files
showing each match with surrounding text, which I read for false positives. That is how I found the Hindi-font
gibberish, ``LLMS'' and ``LL.M''. I also report how many tokens the English filter drops per report --- under 15\% is
expected for English-only reports, over 40\% is flagged for manual checking --- and compare patterns by bank type.}
\end{qlist}

\subsection*{J. Limitations and extensions}
\begin{qlist}
\Q{What is the biggest weakness of your study?}{Identification. Adoption is not random: large private banks adopted
early, while weaker PSU banks adopted late or not at all, and that coincides with the bad-loan cycle, PCA and the mergers.
Fixed effects remove level differences but not different trends, which is exactly what broke parallel trends in the
DiD. Second is measurement: every proxy is indirect. I state both up front.}
\Q{What would you do with six more months?}{Test the mechanisms --- cost-to-income for efficiency and NPAs for credit
quality --- which were planned and deferred. Add a media-tone control so announcements are not confounded with general
news sentiment. Hand-collect deployment dates for the remaining never-treated banks. Run HonestDiD sensitivity bounds.
And extend the panel as FY2026 and later data arrive, since most adoption is recent.}
\Q{What data would you most like to have?}{Bank-level IT and AI spending and technology headcount over time --- a direct
measure of the input rather than disclosure. Second, usage data, for example the share of loans decided by models or
service requests handled by chatbots. Either would replace my indirect proxies.}
\Q{Does this generalise beyond India?}{Partly. The specific numbers are Indian, but the design --- multiple public proxies,
verified adoption dates, staggered DiD with attention to ownership-specific shocks --- applies to other emerging markets
with large state-owned banks, such as China, Brazil or Indonesia. The lesson that ownership-specific shocks can
masquerade as technology effects probably generalises too.}
\Q{What would a regulator take from this?}{That there is no evidence yet that AI announcements or AI talk improve bank
performance, so they should not be read as a signal of strength. Supervision is better focused on what is actually
deployed --- model risk, data governance, fairness in credit decisions --- than on disclosure. And PSU--private
comparisons need to account for the post-clean-up recovery before drawing conclusions about technology.}
\end{qlist}

\subsection*{K. Tools and implementation}
\begin{qlist}
\Q{Why Python for data and R for estimation?}{Each where its tools are strongest. Python is better for downloading,
scraping BSE filings, parsing PDFs and text mining, and assembling the panel. R has the reference implementations of
the estimators I use --- \var{fixest} for fixed effects and Sun--Abraham, \var{did} for Callaway--Sant'Anna --- written
by the methods' authors or maintainers. The two connect through Parquet files.}
\Q{How did you extract text from annual-report PDFs?}{With PyMuPDF. The scripts extract text page by page, normalise
whitespace and hyphenation, cache the extracted text per report so it is not re-parsed, and then count dictionary terms
and English words. Reports yielding fewer than 50 English words per page --- scanned or Hindi-only files --- are set to missing rather than counted as zero.}
\Q{How did you collect announcements and transcripts?}{From BSE's public corporate-announcement endpoints. The scripts
cache every response, can resume after interruption, and stop immediately on authorisation or rate-limit errors (401,
403, 429) rather than hammering the server. Candidate AI announcements were then screened with the dictionary and
verified manually against the two-source rule.}
\Q{How did you query patents?}{Using Google Patents Public Datasets on BigQuery, within the free tier: an SQL query over
publications from 2009 onward, matching harmonised assignee names for all 25 banks, and flagging AI either by CPC class
G06N or by dictionary terms in titles and abstracts. It scanned about 238 GB and returned a single filing.}
\Q{How did you manage versions?}{Git, with the local folder as the source of truth and a public GitHub remote. I committed
after every task with a descriptive message, tagged milestones (\texttt{day1-data}, \texttt{day2-panel},
\texttt{day3-results}), and checked before each commit that no secrets or large raw files were staged. Credentials live
in a git-ignored \texttt{.env} file.}
\end{qlist}

\subsection*{L. About you and the process}
\begin{qlist}
\Q{What was your role vs your supervisor's?}{\emph{(Adjust to what actually happened.)} Prof.\ Sachdeva set the research
question and the overall direction and reviewed the key choices. I wrote the detailed analysis plan, collected and
verified the data, built the text measures and the panel, ran the estimations, and made and logged the day-to-day
decisions --- for example the English-word denominator, dropping the patent proxy and the post-hoc checks. Be specific
about which decisions were yours.}
\Q{What was the hardest decision you made?}{Keeping the strict event-verification rule when it meant falling short of the
70-event target and leaving some banks never-treated --- for example Bank of Maharashtra, whose 2023 machine-learning
model had a bank release but no news report naming it. Loosening the rule would have given more events and a nicer
number, but inconsistent dates would have contaminated the DiD and the event study. I reported the true count, 54.}
\Q{What did you learn?}{That measurement is most of the work: the hardest and most important parts were making the AI
measures mean the same thing across banks, not running regressions. And that pre-specifying the analysis makes it much
easier to report inconvenient results --- the opposite-sign PSU interaction and the DiD that fell apart --- because the
plan, not the result, decides what gets reported.}
\Q{How did you manage a 72-hour sprint?}{By setting daily gates --- clean data and a first set of verified events on day
one; the panel, final events and adoption dates on day two; estimation and robustness on day three --- and by deciding
fallbacks in advance: report the true event count if under target, drop a proxy if its source fails and log it.
Committing after every task meant nothing was lost and progress was always visible.}
\Q{Your CV says ``70+ events''---how many do you have?}{I would correct it straight away: 54 verified events from 88
candidates, 30 used in the event study after exclusions. 70+ was the target; I kept the two-source rule rather than
lower the bar, and the plan said in advance to report the true number. The CV line should say 54.}
\Q{If you started again, what would you do differently?}{Run a feasibility check on each proxy before fixing it in the
spec --- a quick patent search would have shown there were no bank AI patents. Collect transcripts from bank websites
from the start, since four banks only file cover letters on BSE. And start event verification earlier, because it was
the bottleneck for the event count.}
\end{qlist}

\end{document}
